\documentclass[12pt,a4paper]{article}
\usepackage{ctex}
\usepackage{geometry}
\usepackage{graphicx}
\usepackage{listings}
\usepackage{xcolor}
\usepackage{booktabs}
\usepackage{hyperref}

\geometry{left=2.5cm,right=2.5cm,top=2.5cm,bottom=2.5cm}

\definecolor{codebg}{RGB}{245,245,245}
\definecolor{codeframe}{RGB}{200,200,200}

\lstset{
  backgroundcolor=\color{codebg},
  frame=single,
  rulecolor=\color{codeframe},
  basicstyle=\ttfamily\small,
  breaklines=true,
  showstringspaces=false
}

\hypersetup{
  colorlinks=true,
  linkcolor=blue,
  urlcolor=blue
}

\title{Dionysus 用户手册}
\author{FuyuuKu樱}
\date{版本 v0.1.0 \quad 2026-06-16}

\begin{document}

\maketitle
\tableofcontents
\newpage

\section{项目简介}

Dionysus 是一个带有角色陪伴（Companion）功能的 AI Agent 前端。它把聊天会话、Live2D 看板娘、后台角色播报和多个 CLI Agent（kimi / claude / codex / opencode）整合在一个界面里，让你在写代码、查资料的同时，有一个会动、会说的角色陪你。

主要功能：

\begin{itemize}
  \item 多角色切换：内置能天使、凯尔希，也可以自己创建角色。
  \item 语料管理：在线编辑或上传 \texttt{.txt} 语料，塑造角色口吻。
  \item Live2D 模型：选择模型文件夹并一键绑定到角色。
  \item Companion Supervisor：三种后台播报模式可选。
  \item Agent 连接：同时管理多个 CLI Agent 的启用状态与默认模型。
  \item 系统设置：历史记录保留条数、缓存清理、CC Switch 扩展入口。
\end{itemize}

\section{快速启动}

\subsection{启动后端}

\begin{lstlisting}[language=bash]
cd backend
source .venv/bin/activate
python -m uvicorn dionysus_server.main:app --host 0.0.0.0 --port 8765
\end{lstlisting}

后端默认监听 \texttt{http://0.0.0.0:8765}，提供角色、语料、Live2D、Supervisor、Agent 适配器等 REST API。

\subsection{启动前端}

\begin{lstlisting}[language=bash]
cd frontend
npm install   # 首次运行
npm run dev   # 开发模式，默认 http://localhost:5173
\end{lstlisting}

如需生产构建：

\begin{lstlisting}[language=bash]
npm run build
\end{lstlisting}

构建产物位于 \texttt{frontend/dist/}。

\subsection{访问}

打开浏览器访问 \texttt{http://localhost:5173}（开发）或你部署的地址。

\section{角色管理}

\subsection{切换角色}

\begin{enumerate}
  \item 点击左侧导航栏的「角色」图标。
  \item 在「当前角色」下拉框中选择角色。
  \item 右侧 Companion 会立即切换到对应角色，并加载该角色的语料与 Live2D 配置。
\end{enumerate}

\subsection{添加角色}

\begin{enumerate}
  \item 在角色页点击「添加角色」。
  \item 填写：
  \begin{itemize}
    \item \textbf{角色 ID}：英文、数字、下划线，作为文件和目录名。
    \item \textbf{角色名称}：界面上显示的名字。
    \item \textbf{简介}：一句话描述。
  \end{itemize}
  \item 点击「创建」。系统会生成一份默认 YAML，包含系统提示词、Live2D 占位配置、触摸区域、状态表情映射等。
  \item 创建成功后，新角色会出现在下拉框中并自动切换。
\end{enumerate}

\textbf{提示}：ID 创建后不建议修改，因为它会决定运行时配置文件和 Live2D 目录名。

\subsection{编辑语料}

\begin{itemize}
  \item \textbf{方式一}：直接在「语料文本」框中编辑，点击「保存」写入后端。
  \item \textbf{方式二}：点击「选择 .txt 语料」，选中本地文本文件后点击「上传」。上传成功后，文本框会自动刷新。
\end{itemize}

语料会用于构建该角色的系统提示词，影响角色语气、称呼和常用语。

\section{Live2D 模型绑定}

\begin{enumerate}
  \item 准备一个包含 \texttt{.model3.json} 入口文件的 Live2D 模型文件夹。
  \item 在角色页点击「选择模型文件夹」，选中该文件夹。
  \item 按钮旁会显示文件数量，点击「上传并应用」。
  \item 成功后：
  \begin{itemize}
    \item 显示「Live2D 模型已更新：/personas/live2d/<角色ID>/...model3.json」。
    \item Live2D 区域显示「已绑定」和路径。
    \item 右侧 Companion 会尝试加载新模型。
  \end{itemize}
\end{enumerate}

\textbf{注意}：

\begin{itemize}
  \item 首次选择文件夹时，浏览器可能要求用户授权。
  \item 如果右侧模型未立即显示，可点击「重试加载」。
  \item 内置角色（能天使、凯尔希）上传模型时，会复制到运行时目录，不会改动内置文件。
\end{itemize}

\subsection{解绑模型}

点击「解绑模型」可删除已上传的 Live2D 文件并清空角色配置中的 \texttt{model\_path}。

\section{Companion Supervisor}

Supervisor 负责在后台扫描 Agent 会话状态，并让 Companion 以角色口吻播报当前进展。

点击角色页的「后台角色播报」区域，可在三种模式间切换。

\subsection{不接入模型}

完全关闭 Supervisor。Companion 只响应用户主动发送的消息，不会主动播报 Agent 状态。

适合：不需要后台打扰的场景。

\subsection{多开 agent session}

Supervisor 会启动一个独立的 Agent Session（通过你选择的 Adapter ID），把当前会话上下文交给他处理，处理结果由 Companion 播报。

配置项：

\begin{itemize}
  \item \textbf{Adapter ID}：从已启用的 Agent 适配器中选择，例如 \texttt{kimi\_cli}、\texttt{claude\_cli}。
  \item \textbf{扫描间隔（秒）}：Supervisor 轮询间隔，默认 15 秒，最小建议 5 秒。
\end{itemize}

适合：希望用另一个 Agent 实例异步总结或执行任务的场景。

\subsection{DeepSeek API}

Supervisor 直接调用 DeepSeek API 生成播报内容。

配置项：

\begin{itemize}
  \item \textbf{API URL}：默认 \texttt{https://api.deepseek.com/v1/chat/completions}。
  \item \textbf{模型}：默认 \texttt{deepseek-reasoner}，可改为 \texttt{deepseek-chat} 等。
  \item \textbf{API Key}：填写后保存；留空则后端读取环境变量 \texttt{DEEPSEEK\_API\_KEY}。
  \item \textbf{扫描间隔（秒）}：同上。
\end{itemize}

\textbf{安全提示}：前端保存 API Key 时不会明文回显，建议同时通过环境变量配置后端，避免把真实 Key 写入前端代码。

\subsection{保存设置}

修改任意 Supervisor 选项后，点击「保存 Supervisor 设置」。保存成功后按钮旁会显示绿色提示，后端会重新以新配置启动 Supervisor。

\section{Agent 连接设置}

\begin{enumerate}
  \item 点击左侧导航栏的「设置」图标。
  \item 在「Agent 连接」区域可以看到所有已配置的 CLI 适配器：
  \begin{itemize}
    \item \texttt{kimi\_cli}
    \item \texttt{claude\_cli}
    \item \texttt{codex\_cli}
    \item \texttt{opencode\_cli}
  \end{itemize}
  \item 对每个适配器可配置：
  \begin{itemize}
    \item \textbf{启用}：是否允许该适配器被使用。
    \item \textbf{命令路径}：本地 CLI 命令，例如 \texttt{kimi}、\texttt{claude}、\texttt{codex}、\texttt{opencode}。
    \item \textbf{默认模型}：该适配器默认使用的模型名。
  \end{itemize}
  \item 点击「保存并重启」使配置生效。
\end{enumerate}

\textbf{提示}：命令路径必须能在系统 PATH 中找到，否则后端启动 Agent Session 时会失败。

\section{系统设置}

\subsection{历史记录}

\textbf{单会话保留消息数（1-500）}：控制每个会话在本地保存的最大消息数，默认 50。

\subsection{清除本地缓存}

点击「清除本地缓存」可清理浏览器本地存储的临时数据，不会删除后端角色、语料或模型文件。

\subsection{打开 CC Switch}

点击「打开 CC Switch」可跳转或打开 CC Switch 相关扩展/页面。按钮左侧会显示图标。

\section{主题与界面}

Dionysus 使用 Dionysus 设计系统：

\begin{itemize}
  \item 深色玻璃拟态面板
  \item 圆角大按钮、统一边框与阴影
  \item 主题色随角色/配置变化
\end{itemize}

下拉框（\texttt{DionysusSelect}）遵循同样的主题变量，支持亮色与暗色切换。

\section{常见问题}

\begin{enumerate}
  \item \textbf{角色下拉框为空或显示异常？}  
  检查后端是否已启动并可访问 \texttt{/api/personas}；检查浏览器控制台是否有网络错误。

  \item \textbf{语料上传失败？}  
  仅支持 \texttt{.txt} 文件；确保选择了角色后再上传，文件会保存到当前选中角色的语料中。

  \item \textbf{Live2D 模型上传后右侧不显示？}  
  确认文件夹内有 \texttt{.model3.json} 文件；点击 Companion 区域的「重试加载」；检查浏览器控制台是否有资源加载错误。

  \item \textbf{Supervisor 修改后没有生效？}  
  必须点击「保存 Supervisor 设置」；后端日志会显示 \texttt{supervisor\_started mode=<新模式>}，可通过日志确认。

  \item \textbf{Agent 适配器命令找不到？}  
  在系统终端执行对应命令确认已安装；在设置中填写完整路径。
\end{enumerate}

\section{目录与文件速查}

\begin{center}
\begin{tabular}{ll}
\toprule
\textbf{路径} & \textbf{说明} \\
\midrule
\texttt{backend/dionysus\_server/main.py} & 后端主入口 \\
\texttt{backend/config/personas/} & 运行时角色配置、语料、Live2D \\
\texttt{frontend/src/components/Pages/PersonaPage.tsx} & 角色页 \\
\texttt{frontend/src/components/UI/DionysusSelect.tsx} & 统一下拉组件 \\
\texttt{frontend/src/components/Pages/SystemSettingsPage.tsx} & 系统设置页 \\
\bottomrule
\end{tabular}
\end{center}

\end{document}
